\begin{lemma}[Proved: arithmetic defect over a proper curve when radial indices are one]
\label{D:proper-curve-arithmetic-defect}
Let \(k\) be algebraically closed, let \(Y\) be a smooth integral surface
over \(k\), and let \(E\subset Y\) be a smooth proper connected curve
of genus \(g\). Let \(\pi:Y'\to Y\) be a finite dominant generically
separable morphism of degree \(n\), with \(Y'\) normal and every
irreducible component dominating \(Y\).
Assume that every ramification index over the generic point of \(E\)
is one. Write \(\mathcal B=\pi_*\mathcal O_{Y'}\), let \(D_\pi\)
be its trace-discriminant divisor, and put \(C=Y'\times_Y E\).

Then \(C\) is reduced. If its normalization is
\(\coprod_F F^\nu\), with genera \(g_F\), its total normalization
defect
\(\delta(C)=\dim_k H^0(C,\nu_*\mathcal O_{C^\nu}/\mathcal O_C)\)
satisfies
\[
 \delta(C)=\sum_F(1-g_F)-n(1-g)+\frac{D_\pi\cdot E}{2}.
 \tag{AD1}
\]
Factor each finite map \(F^\nu\to E\) into relative Frobenius
followed by a separable map, of separable degree \(s_F\), and
denote its different divisor by \(\Delta_F\). Then
\[
 2\delta(C)+\sum_F\deg\Delta_F
       =2(g-1)\left(n-\sum_Fs_F\right)+D_\pi\cdot E.
 \tag{AD2}
\]
If the discriminants are supported on a fixed finite set of
divisors \(D_j\) and their coefficients satisfy
\(0\leq\operatorname{ord}_{D_j}(D_\pi)\leq c_j n\), then
\(\delta(C)\leq C_E n\), with \(C_E\) independent of the covering.
\end{lemma}

\begin{proof}
A finite dominant normal surface over a regular surface is flat.
Indeed its local rings are Cohen--Macaulay, its fibres have
dimension zero, and miracle flatness applies
\cite[Tag 00R4]{StacksArithmeticDefectB}.
Thus \(\mathcal B\) is a vector bundle of rank \(n\).
At the generic point of \(E\), the assumed indices say that its
base change to \(E\) is a product of fields, although those fields
may be inseparable over \(k(E)\). Hence \(C\) is generically
reduced. Its structure sheaf is torsion-free over \(\mathcal O_E\)
by flatness, so its nilradical, which is generically zero, is zero.

The trace pairing is a map
\(\mathcal B\to\mathcal B^\vee\), generically invertible by
separability. Its determinant is a nonzero section of
\((\det\mathcal B)^{-2}\); by definition its zero divisor is
\(D_\pi\). Therefore
\[
 (\det\mathcal B)^2=\mathcal O_Y(-D_\pi),\qquad
 \deg(\mathcal B|_E)=-\tfrac12 D_\pi\cdot E.
\]
This is the discriminant construction in
\cite[Tag 0BVH]{StacksArithmeticDefectB}.
Finite pushforward and Riemann--Roch on \(E\) give
\[
 \chi(C,\mathcal O_C)
   =\chi(E,\mathcal B|_E)
   =n(1-g)-\tfrac12 D_\pi\cdot E.
\]
Normalization of a reduced proper curve is finite. Its quotient
sheaf is supported at finitely many points, and its Euler
characteristic is its length. The normalization exact sequence
therefore proves (AD1).

The pure-then-separable factorization preserves the genus of each
normalized source component; this is
\cite[Lemmas 53.13.4--53.13.5 and Proposition 53.13.7]
{StacksResidualGenusB}. Riemann--Hurwitz, with the different, gives
\[
 \deg\Delta_F=2g_F-2-s_F(2g-2)
\]
by \cite[Lemma 53.12.2]{StacksResidualGenusB}.
Adding these equalities to twice (AD1) proves (AD2).
Finally there are at most \(n\) components, all their genera are
nonnegative, and
\[
 |D_\pi\cdot E|\leq n\sum_jc_j|D_j\cdot E|.
\]
These estimates in (AD1) prove the claimed linear bound.
\end{proof}

\begin{remark}[Connection with the fixed different bound]
At a divisorial discrete valuation ring, the discriminant
coefficient is \(\sum_F f_F\delta_F\), where \(f_F\) is the full
residue degree and \(\delta_F\) the different exponent.
One can see the formula directly from the trace-dual lattice:
the valuation of the trace determinant is the length of
\(\mathcal B^\vee/\mathcal B\); decomposing over the integral
closure gives length \(f_F\delta_F\) at each extension.
Thus a fixed bound on every \(\delta_F/e_F\), together with
\(\sum_Fe_Ff_F=n\), gives the coefficient bounds in the lemma.

The condition \(e_F=1\) is what makes the Cartier fibre reduced.
When some index is larger, the fibre has nilpotents at its
generic points and its Euler characteristic no longer directly
equals the normalized-genus contribution minus a nonnegative
normalization defect. No bound for that remaining nonreduced
case, and no finite-generation conclusion, is asserted here.
\end{remark}

